\chapter{Chapter05 — Polynomial Threshold Uniform Stability}
%%% generated by the blueprint pipeline from FABL.Chapter05.PolynomialThresholdUniformStability
\section{Overview}

Declarations of the Lean module \texttt{FABL.Chapter05.PolynomialThresholdUniformStability}.

\section{Declarations}

\begin{definition}[Polynomial threshold class]
\lean{FABL.polynomialThresholdClass}
\statementsource{Analysis-of-Boolean-Functions-by-Ryan-ODonnell}{pdf-ffbe9c7b3306-p133-b003}
\label{FABL.polynomialThresholdClass}
\leanok
\uses{FABL.BooleanClass, FABL.BooleanFunction, FABL.IsPolynomialThreshold}
For each natural number $k$, define a Boolean function class by declaring that, for
every natural number $n$, its members are exactly the Boolean functions on $n$
coordinates that are polynomial thresholds of degree at most $k$.
\end{definition}

\begin{definition}[Identified frequency set]
\lean{FABL.identifiedFrequency}
\label{FABL.identifiedFrequency}
\leanok
Let $n,m$ be nonnegative integers, let $\pi:\{0,\ldots,n-1\}\to\{0,\ldots,m-1\}$, and
let $S\subseteq\{0,\ldots,n-1\}$. The identified frequency set consists of those $j$ for
which the number of elements of $S$ mapped to $j$ by $\pi$ is odd.
\end{definition}

\begin{theorem}[Bound on identified frequency size]
\lean{FABL.identifiedFrequency_card_le}
\label{FABL.identifiedFrequency_card_le}
\leanok
\uses{FABL.identifiedFrequency}
\difficulty{4}
Let $n,m\ge0$, let $\pi\colon\{1,\ldots,n\}\to\{1,\ldots,m\}$ be a map, and let
$S\subseteq\{1,\ldots,n\}$. Form the set of those target indices $j$ for which the
number of elements of $S$ mapped to $j$ is odd. The cardinality of this set is at most
$|S|$.
\end{theorem}

\begin{theorem}[Product of signs over a fiber]
\lean{FABL.fiberSignProduct_eq_ite_odd}
\label{FABL.fiberSignProduct_eq_ite_odd}
\leanok
\uses{FABL.signValue, FABL.signValue_eq_neg_one_or_one}
\difficulty{4}
Let $n,m$ be natural numbers, let $S\subseteq\{0,\ldots,n-1\}$, let
$\pi:\{0,\ldots,n-1\}\to\{0,\ldots,m-1\}$, let $w$ be an $m$-tuple of signs, and choose
$j\in\{0,\ldots,m-1\}$. The product of the sign at $j$ over all $i\in S$ with $\pi(i)=j$
equals that sign if the number of such indices is odd, and equals $1$ otherwise.
\end{theorem}

\begin{theorem}[Monomial on a partitioned input]
\lean{FABL.monomial_partitionInput}
\label{FABL.monomial_partitionInput}
\leanok
\uses{FABL.fiberSignProduct_eq_ite_odd, FABL.identifiedFrequency, FABL.monomial, FABL.signValue}
\difficulty{5}
Let $n$ and $m$ be natural numbers, let $S$ be a subset of the $n$ input coordinates,
let $z$ be a vector of $n$ signs, let $\pi$ map the $n$ input coordinates to the $m$
output coordinates, and let $w$ be a vector of $m$ signs. Form $R$ from those output
coordinates that are the image under $\pi$ of an odd number of elements of $S$. The
monomial indexed by $S$ at the vector with coordinates $z_iw_{\pi(i)}$ equals the
monomial indexed by $S$ at $z$ times the monomial indexed by $R$ at $w$.
\end{theorem}

\begin{definition}[Partition polynomial coefficient]
\lean{FABL.partitionPolynomialCoefficient}
\label{FABL.partitionPolynomialCoefficient}
\leanok
\uses{FABL.fourierCoeff, FABL.identifiedFrequency, FABL.monomial}
Let $p:\{−1,1\}^n\to\mathbb R$, $z\in\{−1,1\}^n$, and
$\pi:\{0,\ldots,n-1\}\to\{0,\ldots,m-1\}$, and let $T\subseteq\{0,\ldots,m-1\}$. The
partition polynomial coefficient at $T$ is $\sum_{S:\,J_\pi(S)=T}\widehat
p(S)\,\chi_S(z)$, where $J_\pi(S)$ is the set of indices receiving an odd number of
elements of $S$, $\widehat p(S)$ is the Fourier coefficient of $p$ at $S$, and
$\chi_S(z)=\prod_{i\in S}z_i$.
\end{definition}

\begin{definition}[Partition polynomial]
\lean{FABL.partitionPolynomial}
\label{FABL.partitionPolynomial}
\leanok
\uses{FABL.multilinearPolynomial, FABL.partitionPolynomialCoefficient}
Let $p:\{−1,1\}^n\to\mathbb R$, $z\in\{−1,1\}^n$, and
$\pi:\{0,\ldots,n-1\}\to\{0,\ldots,m-1\}$. The partition polynomial is the multilinear
polynomial on $\{−1,1\}^m$ whose coefficient at each set $T\subseteq\{0,\ldots,m-1\}$ is
$\sum_{S:\,J_\pi(S)=T}\widehat p(S)\,\chi_S(z)$, with $J_\pi(S)$ the odd-multiplicity
image set under $\pi$.
\end{definition}

\begin{theorem}[Evaluation of the partition polynomial]
\lean{FABL.partitionPolynomial_apply}
\label{FABL.partitionPolynomial_apply}
\leanok
\uses{FABL.fourierCoeff, FABL.fourier_expansion, FABL.identifiedFrequency, FABL.monomial, FABL.monomial_partitionInput, FABL.multilinearPolynomial, FABL.partitionPolynomial, FABL.partitionPolynomialCoefficient}
\difficulty{5}
Let $p$ be a real-valued function on the $n$-dimensional sign cube, let $z$ be an
$n$-dimensional sign vector, and let $\pi$ map the $n$ coordinates to $m$ coordinates.
The partition polynomial associated with $p,z,\pi$, evaluated at any $m$-dimensional
sign vector $w$, equals $p\bigl((z_iw_{\pi(i)})_{i=1}^n\bigr)$.
\end{theorem}

\begin{theorem}[Degree of a partition polynomial]
\lean{FABL.fourierDegree_partitionPolynomial_le}
\label{FABL.fourierDegree_partitionPolynomial_le}
\leanok
\uses{FABL.fourierCoeff, FABL.fourierDegree, FABL.fourierDegree_le_iff, FABL.fourier_expansion_unique, FABL.identifiedFrequency, FABL.identifiedFrequency_card_le, FABL.partitionPolynomial, FABL.partitionPolynomialCoefficient}
\difficulty{5}
Let $p$ be a real-valued function on the Boolean cube with Fourier degree at most $k$,
let $z$ be a point of that cube, and let $\\pi:[n]\\to[m]$ be any map. The partition
polynomial formed from $p$, $z$, and $\\pi$ has Fourier degree at most $k$.
\end{theorem}

\begin{theorem}[Thresholds under coordinate partitioning]
\lean{FABL.isPolynomialThreshold_partitionFunction}
\label{FABL.isPolynomialThreshold_partitionFunction}
\leanok
\uses{FABL.BooleanFunction, FABL.IsPolynomialThreshold, FABL.fourierDegree_partitionPolynomial_le, FABL.partitionFunction, FABL.partitionPolynomial, FABL.partitionPolynomial_apply, FABL.thresholdSign}
\difficulty{3}
Let $n,m,k$ be nonnegative integers, let $f:\{-1,1\}^n\to\{-1,1\}$ have a polynomial
threshold representation of Fourier degree at most $k$, let $z\in\{-1,1\}^n$, and let
$\pi:\{1,\ldots,n\}\to\{1,\ldots,m\}$ be any map. The function $w\mapsto
f\bigl((z_iw_{\pi(i)})_{i=1}^n\bigr)$ on $\{-1,1\}^m$ also has a polynomial threshold
representation of Fourier degree at most $k$.
\end{theorem}

\begin{theorem}[Closure under negating input variables]
\lean{FABL.polynomialThresholdClass_closedUnderNegatingInputVariables}
\label{FABL.polynomialThresholdClass_closedUnderNegatingInputVariables}
\leanok
\uses{FABL.IsClosedUnderNegatingInputVariables, FABL.IsPolynomialThreshold, FABL.isPolynomialThreshold_partitionFunction, FABL.negateInputVariables, FABL.partitionFunction, FABL.polynomialThresholdClass}
\difficulty{3}
For every natural number $k$, the class of Boolean functions admitting a polynomial
threshold representation of Fourier degree at most $k$ is closed under negating input
variables: for any number of input coordinates, any assignment of signs to those
coordinates, and any function in the class, the function obtained by multiplying each
input coordinate by its assigned sign belongs to the same class.
\end{theorem}

\begin{theorem}[Closure under identifying input variables]
\lean{FABL.polynomialThresholdClass_closedUnderIdentifyingInputVariables}
\label{FABL.polynomialThresholdClass_closedUnderIdentifyingInputVariables}
\leanok
\uses{FABL.IsClosedUnderIdentifyingInputVariables, FABL.IsPolynomialThreshold, FABL.identifyInputVariables, FABL.isPolynomialThreshold_partitionFunction, FABL.partitionFunction, FABL.polynomialThresholdClass}
\difficulty{3}
For every natural number $k$, the class of Boolean functions admitting a polynomial
threshold representation of Fourier degree at most $k$ is closed under identifying input
variables: for any natural numbers $n,m$, any map from the $n$ input coordinates to the
$m$ input coordinates, and any function in the class on $n$ coordinates, the function
obtained by substituting the coordinate indexed by the map for each original input
coordinate belongs to the class on $m$ coordinates.
\end{theorem}

\begin{definition}[Uniformly sublinear total influence]
\lean{FABL.HasUniformlySublinearTotalInfluence}
\statementsource{Analysis-of-Boolean-Functions-by-Ryan-ODonnell}{pdf-ffbe9c7b3306-p139-b003}
\label{FABL.HasUniformlySublinearTotalInfluence}
\leanok
\uses{FABL.BooleanClass, FABL.BooleanFunction, FABL.BooleanFunction.toReal, FABL.UnitProbability, FABL.totalInfluence}
A Boolean function class has uniformly sublinear total influence if there is a sequence
$\eta_n\in[0,1]$ tending to zero such that, for every natural number $n$ and every
function $f$ in the class at dimension $n$, the total influence of its real-valued
version is at most $n\eta_n$.
\end{definition}

\begin{definition}[Uniform square-root noise sensitivity bound]
\lean{FABL.HasUniformPolynomialThresholdNoiseSensitivitySqrtBound}
\statementsource{Analysis-of-Boolean-Functions-by-Ryan-ODonnell}{pdf-ffbe9c7b3306-p139-b003}
\label{FABL.HasUniformPolynomialThresholdNoiseSensitivitySqrtBound}
\leanok
\uses{FABL.BooleanFunction, FABL.IsPolynomialThreshold, FABL.PositiveHalfNoiseParameter, FABL.noiseSensitivity}
There is a constant $C\geq0$ such that, for every natural number $k$, every natural
number $n$, every Boolean function $f$ on $n$ coordinates that is a polynomial threshold
of degree at most $k$, and every real noise parameter $\delta$ with
$0<\delta\leq\tfrac12$, the noise sensitivity of $f$ at parameter $\delta$ is at most
$Ck\sqrt{\delta}$.
\end{definition}

\begin{definition}[Uniform influence bound for polynomial thresholds]
\lean{FABL.HasUniformPolynomialThresholdTotalInfluenceSqrtBound}
\statementsource{Analysis-of-Boolean-Functions-by-Ryan-ODonnell}{pdf-ffbe9c7b3306-p139-b003}
\label{FABL.HasUniformPolynomialThresholdTotalInfluenceSqrtBound}
\leanok
\uses{FABL.BooleanFunction, FABL.BooleanFunction.toReal, FABL.IsPolynomialThreshold, FABL.totalInfluence}
There is a constant $C\ge 0$ such that, for every natural number $k$, every natural
number $n$, and every Boolean function $f$ on the $n$-dimensional Boolean cube that is a
polynomial threshold function of degree at most $k$, the total influence of its
associated real-valued function is at most $Ck\sqrt n$.
\end{definition}

\begin{theorem}[Influence bound from noise sensitivity]
\lean{FABL.totalInfluence_le_constant_mul_dimension_of_noiseSensitivity_inverseDimension_le}
\proofsource{Analysis-of-Boolean-Functions-by-Ryan-ODonnell}{
  pdf-ffbe9c7b3306-p064-b009,
  pdf-ffbe9c7b3306-p139-b003
}
\label{FABL.totalInfluence_le_constant_mul_dimension_of_noiseSensitivity_inverseDimension_le}
\leanok
\uses{FABL.BooleanFunction, FABL.BooleanFunction.toReal, FABL.averageInfluence, FABL.averageInfluence_mul_one_sub_exp_neg_two_div_two_le_noiseSensitivity, FABL.noiseSensitivity, FABL.totalInfluence}
\difficulty{4}
Let $n\ge 2$, let $f$ be a Boolean function on the $n$-dimensional cube, and let
$E\in\mathbb R$. If the probability that $f$ changes under noise that independently
flips each coordinate with probability $1/n$ is at most $E$, then the total influence of
the real-valued sign representation of $f$ is at most $\displaystyle
\frac{2}{1-e^{-2}}\,nE$.
\end{theorem}

\begin{theorem}[Influence bound from noise sensitivity]
\lean{FABL.totalInfluence_le_two_div_one_sub_exp_neg_two_mul_dimension_mul_noiseModulus}
\label{FABL.totalInfluence_le_two_div_one_sub_exp_neg_two_mul_dimension_mul_noiseModulus}
\leanok
\uses{FABL.BooleanFunction, FABL.BooleanFunction.toReal, FABL.HalfNoiseParameter, FABL.IsPolynomialThreshold, FABL.UnitProbability, FABL.noiseSensitivity, FABL.totalInfluence, FABL.totalInfluence_le_constant_mul_dimension_of_noiseSensitivity_inverseDimension_le}
\difficulty{3}
Let $n,m,k$ be natural numbers with $2\le n$. Suppose that for every natural number $r$,
every Boolean function on the $r$-dimensional sign cube of polynomial threshold degree
at most $k$, and every noise parameter $\delta\in[0,1/2]$, the noise sensitivity at
$\delta$ is at most $\varepsilon(\delta)$, where $\varepsilon$ takes values in $[0,1]$.
If $f$ is a Boolean function on the $n$-dimensional sign cube of polynomial threshold
degree at most $k$, then its total influence is at most
$\frac{2n}{1-e^{-2}}\,\varepsilon(1/n)$.
\end{theorem}

\begin{definition}[Inverse-dimension half-noise parameter]
\lean{FABL.inverseDimensionHalfNoiseParameter}
\label{FABL.inverseDimensionHalfNoiseParameter}
\leanok
\uses{FABL.HalfNoiseParameter}
For each nonnegative integer $n$, define a parameter in $[0,1/2]$ to be $1/n$ when
$n\ge2$, and $1/2$ when $n<2$.
\end{definition}

\begin{theorem}[Inverse dimension parameter tends to zero]
\lean{FABL.tendsto_inverseDimensionHalfNoiseParameter}
\label{FABL.tendsto_inverseDimensionHalfNoiseParameter}
\leanok
\uses{FABL.HalfNoiseParameter, FABL.inverseDimensionHalfNoiseParameter}
\difficulty{4}
For each natural number $n$, define $\delta_n=1/n$ when $n\ge2$, and $\delta_n=1/2$
otherwise. As $n$ tends to infinity, $\delta_n$ tends to $0$ through positive values in
$[0,1/2]$.
\end{theorem}

\begin{definition}[Influence modulus from noise modulus]
\lean{FABL.influenceModulusOfNoiseModulus}
\label{FABL.influenceModulusOfNoiseModulus}
\leanok
\uses{FABL.HalfNoiseParameter, FABL.UnitProbability, FABL.inverseDimensionHalfNoiseParameter}
Let $\varepsilon:[0,1/2]\to[0,1]$ and let $n$ be a nonnegative integer. If $n\ge2$,
define the resulting parameter to be
$\min\!\left(1,\frac{2}{1-e^{-2}}\varepsilon(1/n)\right)$; if $n<2$, define it to be
$1$.
\end{definition}

\begin{theorem}[Vanishing influence modulus]
\lean{FABL.tendsto_influenceModulusOfNoiseModulus}
\label{FABL.tendsto_influenceModulusOfNoiseModulus}
\leanok
\uses{FABL.HalfNoiseParameter, FABL.UnitProbability, FABL.influenceModulusOfNoiseModulus, FABL.inverseDimensionHalfNoiseParameter, FABL.tendsto_inverseDimensionHalfNoiseParameter}
\difficulty{4}
Let $\varepsilon$ assign to each $\delta\in[0,1/2]$ a number in $[0,1]$. Suppose
$\varepsilon(\delta)$ tends to $0$ as positive $\delta$ tends to $0$ within $[0,1/2]$.
For each natural number $n$, define the influence modulus to be $1$ if $n<2$, and
otherwise to be $\min\{1,\frac{2}{1-e^{-2}}\varepsilon(1/n)\}$. This modulus tends to
$0$ as $n$ tends to infinity.
\end{theorem}

\begin{definition}[Inverse noise floor at zero]
\lean{FABL.inverseNoiseFloorAtZero}
\label{FABL.inverseNoiseFloorAtZero}
\leanok
\uses{FABL.HalfNoiseParameter, FABL.PositiveHalfNoiseParameter, FABL.inverseNoiseFloor}
Let $\delta\in[0,1/2]$. If $\delta>0$, the inverse noise floor at $\delta$ is the
inverse noise floor evaluated at $\delta$ as a positive parameter; if $\delta=0$, it is
$1$.
\end{definition}

\begin{theorem}[Divergence of the inverse noise floor]
\lean{FABL.tendsto_inverseNoiseFloorAtZero}
\label{FABL.tendsto_inverseNoiseFloorAtZero}
\leanok
\uses{FABL.HalfNoiseParameter, FABL.inverseNoiseFloorAtZero}
\difficulty{5}
For $0\leq\delta\leq\tfrac12$, let $N(\delta)=\lfloor 1/\delta\rfloor$ when $\delta>0$
and $N(0)=1$. As $\delta$ tends to $0$ through positive values in $[0,\tfrac12]$,
$N(\delta)$ tends to infinity.
\end{theorem}

\begin{definition}[Noise modulus from influence modulus]
\lean{FABL.noiseModulusOfInfluenceModulus}
\label{FABL.noiseModulusOfInfluenceModulus}
\leanok
\uses{FABL.HalfNoiseParameter, FABL.PositiveHalfNoiseParameter, FABL.UnitProbability, FABL.inverseNoiseFloor}
Let $\eta:\mathbb N\to[0,1]$ and let $\delta\in[0,1/2]$. If $\delta>0$, define the
resulting parameter to be $\eta$ evaluated at the inverse noise floor of $\delta$; if
$\delta=0$, define it to be $0$.
\end{definition}

\begin{theorem}[Continuity of the induced noise modulus]
\lean{FABL.tendsto_noiseModulusOfInfluenceModulus}
\label{FABL.tendsto_noiseModulusOfInfluenceModulus}
\leanok
\uses{FABL.HalfNoiseParameter, FABL.UnitProbability, FABL.inverseNoiseFloorAtZero, FABL.noiseModulusOfInfluenceModulus, FABL.tendsto_inverseNoiseFloorAtZero}
\difficulty{3}
Let $\eta(n)$ be a number in $[0,1]$ for each natural number $n$, and suppose
$\eta(n)\to 0$ as $n\to\infty$. Define a function on $[0,1/2]$ by taking its value at
$\delta>0$ to be $\eta(\lfloor1/\delta\rfloor)$, and its value at $\delta=0$ to be $0$.
This function tends to $0$ as $\delta\to0$ through positive values in $[0,1/2]$.
\end{theorem}

\begin{theorem}[Noise stability and sublinear total influence]
\lean{FABL.polynomialThresholdClass_uniformlyNoiseStable_iff_uniformlySublinearTotalInfluence}
\label{FABL.polynomialThresholdClass_uniformlyNoiseStable_iff_uniformlySublinearTotalInfluence}
\leanok
\uses{FABL.BooleanFunction.toReal, FABL.HalfNoiseParameter, FABL.HasUniformlySublinearTotalInfluence, FABL.IsPolynomialThreshold, FABL.IsUniformlyNoiseStable, FABL.PositiveHalfNoiseParameter, FABL.influenceModulusOfNoiseModulus, FABL.inverseDimensionHalfNoiseParameter, FABL.inverseNoiseFloor, FABL.noiseModulusOfInfluenceModulus, FABL.noiseSensitivity, FABL.noiseSensitivity_eq_sum_level, FABL.noiseSensitivity_le_inverseNoiseFloor_totalInfluenceBound, FABL.polynomialThresholdClass, FABL.polynomialThresholdClass_closedUnderIdentifyingInputVariables, FABL.polynomialThresholdClass_closedUnderNegatingInputVariables, FABL.tendsto_influenceModulusOfNoiseModulus, FABL.tendsto_noiseModulusOfInfluenceModulus, FABL.totalInfluence, FABL.totalInfluence_le_two_div_one_sub_exp_neg_two_mul_dimension_mul_noiseModulus, FABL.totalInfluence_toReal_mem_Icc}
\difficulty{6}
For every natural number $k$, the class of Boolean functions admitting a polynomial
threshold representation of Fourier degree at most $k$ is uniformly noise stable if and
only if it has uniformly sublinear total influence. Here uniform noise stability means
there is a function assigning a probability in $[0,1]$ to each noise parameter in
$[0,1/2]$, tending to zero as the parameter tends to zero through positive values, that
bounds the noise sensitivity of every function in the class at every such parameter.
Uniformly sublinear total influence means there is a sequence of probabilities in
$[0,1]$ tending to zero that bounds, for every dimension $n$, the total influence of
each function in the class by $n$ times the corresponding probability.
\end{theorem}

\begin{theorem}[Influence and noise sensitivity in dimension one]
\lean{FABL.totalInfluence_dimension_one_eq_two_mul_noiseSensitivity_half}
\label{FABL.totalInfluence_dimension_one_eq_two_mul_noiseSensitivity_half}
\leanok
\uses{FABL.BooleanFunction, FABL.BooleanFunction.toReal, FABL.noiseSensitivity, FABL.noiseSensitivity_eq_sum_level, FABL.totalInfluence, FABL.totalInfluence_eq_sum_level_mul_fourierWeight}
\difficulty{3}
For every Boolean function $f$ on the one-dimensional cube, its total influence equals
twice its noise sensitivity at noise rate $1/2$.
\end{theorem}

\begin{theorem}[Uniform polynomial threshold noise sensitivity sqrt bound iff total influence sqrt bound]
\lean{FABL.uniformPolynomialThreshold_noiseSensitivitySqrtBound_iff_totalInfluenceSqrtBound}
\statementsource{Analysis-of-Boolean-Functions-by-Ryan-ODonnell}{pdf-ffbe9c7b3306-p139-b003}
\label{FABL.uniformPolynomialThreshold_noiseSensitivitySqrtBound_iff_totalInfluenceSqrtBound}
\leanok
\uses{FABL.BooleanFunction.toReal, FABL.HasUniformPolynomialThresholdNoiseSensitivitySqrtBound, FABL.HasUniformPolynomialThresholdTotalInfluenceSqrtBound, FABL.PositiveHalfNoiseParameter, FABL.inverseNoiseFloor, FABL.noiseSensitivity, FABL.noiseSensitivity_le_inverseNoiseFloor_totalInfluenceBound, FABL.polynomialThresholdClass, FABL.polynomialThresholdClass_closedUnderIdentifyingInputVariables, FABL.polynomialThresholdClass_closedUnderNegatingInputVariables, FABL.sqrt_inverseNoiseFloor_le_sqrt_three_halves_mul_sqrt, FABL.totalInfluence, FABL.totalInfluence_dimension_one_eq_two_mul_noiseSensitivity_half, FABL.totalInfluence_le_constant_mul_dimension_of_noiseSensitivity_inverseDimension_le, FABL.totalInfluence_toReal_mem_Icc}
\difficulty{6}
% dataset-pipeline: stub -- prose pending build_informal_statements.py
Statement pending regeneration.
\end{theorem}
